%%%%%%%%%%%%%%%%%%%%%%%%%%%%%%%%%%%%%%%%%%%%%%%%%%%%%%
\xchapter{Desenvolvimento do FitWorkUp}{}
\label{cap:fitworkup}
%%%%%%%%%%%%%%%%%%%%%%%%%%%%%%%%%%%%%%%%%%%%%%%%%%%%%%

Este capítulo descreve o artefato desenvolvido, suas funcionalidades, as tecnologias empregadas, a arquitetura adotada e os principais elementos de modelagem. A descrição foi elaborada a partir da implementação do aplicativo Android e da API, de modo que recursos ainda não concluídos sejam tratados como limitações ou trabalhos futuros.

\section{Visão Geral da Solução}
\label{sec:fitworkup_visao_geral}

O \textit{FitWorkUp} é composto por um aplicativo Android, uma API REST e um banco de dados relacional. O aplicativo permite registrar caminhadas e corridas, acompanhar métricas, consultar o histórico, visualizar a progressão do perfil e acessar recursos sociais e de gamificação. A API centraliza autenticação, regras de negócio, persistência, validação das atividades e concessão de recompensas.

Durante uma atividade, o cliente coleta informações de localização e movimento. Ao término, envia à API dados como distância, quantidade de passos, velocidade média e indicadores produzidos pela análise local. O servidor executa verificações adicionais e classifica o registro como aprovado, em análise ou rejeitado. Somente atividades aprovadas geram experiência e \textit{FitCoins}. Esse processo busca reduzir registros inconsistentes, mas não garante a detecção de toda forma de manipulação.

\section{Funcionalidades Implementadas}
\label{sec:fitworkup_funcionalidades}

As funcionalidades foram agrupadas de acordo com a responsabilidade de cada módulo.

\subsection{Autenticação e Sessão}

O módulo de autenticação permite cadastrar usuários, autenticar credenciais, manter uma sessão por token JWT e solicitar recuperação de senha. O aplicativo armazena as informações necessárias à sessão com o \textit{DataStore}; nas rotas protegidas, o token é acrescentado às requisições pelo cliente HTTP. A API utiliza \textit{Spring Security} para validar o token e identificar o usuário autenticado.

\subsection{Registro de Atividades}

O usuário pode iniciar uma atividade individual ou participar de uma sessão em grupo. Durante o treino, a interface apresenta métricas como duração, distância, passos e percurso no mapa. O \texttt{WorkoutSensorService} mantém a coleta enquanto a atividade está em andamento, inclusive quando a tela deixa de ser o elemento visível do aplicativo.

A análise local produz passos aceitos, passos retidos, pontuação de risco e justificativas de suspeita. No servidor, o \texttt{FraudDetectionService} verifica a presença dos dados obrigatórios, a plausibilidade da velocidade e a relação entre passos e distância. Dependendo do resultado, o registro recebe o estado aprovado, em análise ou rejeitado.

\subsection{Perfil, Progresso e Conquistas}

O perfil reúne nível, experiência, sequência de uso, distância acumulada, avatar, título e conquistas. Atividades aprovadas concedem experiência e moedas virtuais conforme as regras implementadas no \texttt{ActivityService}. O \texttt{GamificationService} aplica multiplicadores ativos e atualiza o nível quando a experiência acumulada alcança o limite definido para a progressão.

\subsection{Ranking e Recursos Sociais}

O ranking semanal apresenta usuários ordenados pelo desempenho considerado pela API. O módulo social permite pesquisar perfis, enviar e responder solicitações de amizade, consultar amigos e visualizar o perfil público de usuários com vínculo aceito. Também existem sessões de atividade em grupo, com criação de sala, entrada por código, indicação de prontidão e início pelo anfitrião.

\subsection{Loja, Inventário e Efeitos Temporários}

A loja apresenta itens cadastrados no servidor e permite adquiri-los com \textit{FitCoins}. O inventário registra quantidade e estado de equipamento. Alguns itens produzem efeitos temporários, como multiplicadores de experiência ou moedas, representados no servidor pela entidade \texttt{UserBoost}. A compra verifica o saldo e é processada na camada de serviço.

\subsection{Resumo Diário e Mensagens de Incentivo}

O resumo diário agrega passos, distância e estimativa de calorias das atividades válidas. O serviço atual de mensagens de incentivo utiliza regras baseadas na quantidade diária de passos. Embora versões anteriores da documentação mencionassem \textit{Spring AI} e Google Gemini, essa integração não está presente nas dependências atuais; portanto, não é descrita como funcionalidade implementada.

\section{Tecnologias Utilizadas}
\label{sec:fitworkup_tecnologias}

A Tabela~\ref{tab:tecnologias_fitworkup} apresenta as tecnologias identificadas na implementação e o papel de cada uma.

\begin{longtable}{p{0.20\textwidth}p{0.25\textwidth}p{0.46\textwidth}}
\caption{Tecnologias utilizadas no FitWorkUp}\label{tab:tecnologias_fitworkup}\\
\toprule
\textbf{Camada} & \textbf{Tecnologia} & \textbf{Uso no projeto} \\
\midrule
\endfirsthead
\toprule
\textbf{Camada} & \textbf{Tecnologia} & \textbf{Uso no projeto} \\
\midrule
\endhead
\bottomrule
\endfoot
Aplicativo & Kotlin e Android SDK & Implementação do cliente Android, serviços, acesso aos sensores e regras executadas no dispositivo. \\
Interface & Jetpack Compose e Material 3 & Construção declarativa das telas, componentes e estados visuais. \\
Arquitetura móvel & ViewModel, Hilt e Navigation Compose & Gerenciamento de estado, injeção de dependências e navegação entre telas. \\
Persistência local & Room e DataStore & Armazenamento local de atividades e preferências relacionadas à sessão. \\
Processamento em segundo plano & Foreground Service e WorkManager & Coleta durante o treino e execução de tarefas de sincronização. \\
Localização e mapas & Fused Location Provider e Maps Compose & Obtenção de localização e apresentação do percurso. \\
Comunicação & Retrofit, OkHttp e Gson & Consumo da API REST, interceptação das requisições e conversão de JSON. \\
Servidor & Java 21 e Spring Boot 3.3.5 & Implementação da API e inicialização dos componentes do servidor. \\
Segurança & Spring Security, JWT e BCrypt & Autenticação, autorização de rotas e armazenamento protegido das senhas. \\
Persistência remota & Spring Data JPA e PostgreSQL & Mapeamento objeto-relacional e armazenamento dos dados da aplicação. \\
Documentação da API & Springdoc OpenAPI & Exposição da especificação e da interface de consulta dos endpoints. \\
\end{longtable}

\section{Arquitetura do Sistema}
\label{sec:fitworkup_arquitetura}

O sistema segue uma arquitetura cliente-servidor. No aplicativo, a interface em Compose observa estados mantidos pelos \textit{ViewModels}. As abstrações de repositório separam as regras usadas pela apresentação dos detalhes de acesso à API e à base local. No servidor, controladores recebem as requisições HTTP e delegam o processamento a serviços, que utilizam repositórios JPA para acessar o PostgreSQL.

\begin{figure}[H]
    \centering
    \resizebox{\textwidth}{!}{\begin{tikzpicture}[
    font=\small,
    node distance=0.9cm and 1.2cm,
    bloco/.style={draw, rounded corners, align=center, minimum height=1.05cm, minimum width=3.3cm, fill=blue!6},
    externo/.style={draw, dashed, rounded corners, align=center, minimum height=0.9cm, minimum width=3.0cm, fill=gray!8},
    fluxo/.style={->, >=stealth, thick}
]
    \node[bloco] (ui) {Aplicativo Android\\Kotlin + Jetpack Compose};
    \node[bloco, below=of ui] (vm) {Apresentação e estado\\Telas + ViewModels};
    \node[bloco, below=of vm] (domain) {Domínio\\Casos de uso + Repositórios};
    \node[bloco, below left=of domain] (local) {Persistência local\\Room + DataStore};
    \node[bloco, below right=of domain] (sensors) {Coleta de atividade\\GPS + sensores de movimento};

    \node[bloco, right=3.0cm of vm] (api) {API REST\\Spring Boot + Spring Security};
    \node[bloco, below=of api] (services) {Serviços de negócio\\atividade, gamificação, loja,\\ranking, amizades e grupos};
    \node[bloco, below=of services] (jpa) {Persistência no servidor\\Spring Data JPA};
    \node[bloco, below=of jpa] (db) {PostgreSQL};

    \node[externo, above=of api] (maps) {Google Maps e serviços\\de localização};
    \node[externo, right=of services] (mail) {Servidor de e-mail\\recuperação de senha};

    \draw[fluxo] (ui) -- (vm);
    \draw[fluxo] (vm) -- (domain);
    \draw[fluxo] (domain) -- (local);
    \draw[fluxo] (domain) -- (sensors);
    \draw[fluxo,<->] (vm) -- node[above, align=center]{HTTPS/JSON\\JWT} (api);
    \draw[fluxo] (api) -- (services);
    \draw[fluxo] (services) -- (jpa);
    \draw[fluxo] (jpa) -- (db);
    \draw[fluxo,<->] (sensors) -- (maps);
    \draw[fluxo,<->] (services) -- (mail);
\end{tikzpicture}
}
    \caption{Arquitetura do ecossistema FitWorkUp}
    \label{fig:arquitetura_fitworkup}
\end{figure}

A comunicação entre o aplicativo e a API ocorre por HTTP com dados em JSON. Nas operações autenticadas, o JWT identifica a sessão. A coleta de sensores permanece no dispositivo; apenas as métricas e os indicadores necessários ao registro são enviados ao servidor. Serviços externos são empregados para localização, mapas e envio de mensagens de recuperação de senha.

\section{Modelagem de Classes}
\label{sec:fitworkup_classes}

A Figura~\ref{fig:classes_fitworkup} apresenta um recorte das classes envolvidas no registro de uma atividade. O recorte foi escolhido por representar o fluxo técnico central do sistema e por evidenciar a divisão entre cliente e servidor.

\begin{figure}[H]
    \centering
    \resizebox{\textwidth}{!}{\begin{tikzpicture}[
    font=\scriptsize,
    node distance=0.7cm and 0.9cm,
    classe/.style={draw, rounded corners, align=left, text width=3.2cm, inner sep=5pt, fill=blue!5},
    interface/.style={draw, dashed, rounded corners, align=left, text width=3.2cm, inner sep=5pt, fill=gray!7},
    dep/.style={->, >=stealth, thick},
    impl/.style={->, >=stealth, dashed, thick}
]
    \node[classe] (screen) {\textbf{WorkoutScreen}\\exibe estado do treino\\envia ações do usuário};
    \node[classe, below=of screen] (vm) {\textbf{WorkoutViewModel}\\coordena o estado\\inicia e finaliza atividade};
    \node[interface, below left=of vm] (repoi) {\textit{ActivityRepository}\\registrar atividade\\consultar histórico};
    \node[classe, below=of repoi] (repo) {\textbf{ActivityRepositoryImpl}\\integra API e base local};
    \node[classe, below right=of vm] (sensor) {\textbf{WorkoutSensorService}\\coleta localização e movimento\\mantém execução em primeiro plano};
    \node[classe, below=of sensor] (evaluator) {\textbf{StepAntiFraudEvaluator}\\avalia passos aceitos/retidos\\produz risco e justificativas};

    \node[classe, right=3.0cm of vm] (controller) {\textbf{ActivityController}\\recebe e valida requisições\\retorna o estado da atividade};
    \node[classe, below=of controller] (activityservice) {\textbf{ActivityService}\\orquestra persistência, validação,\\recompensas e conquistas};
    \node[classe, below left=of activityservice] (fraud) {\textbf{FraudDetectionService}\\classifica a atividade em\\aprovada, revisão ou rejeitada};
    \node[classe, below right=of activityservice] (game) {\textbf{GamificationService}\\atribui XP e FitCoins\\aplica multiplicadores e níveis};
    \node[interface, below=of activityservice] (jparepo) {\textit{ActivityRepository (JPA)}\\save, consultas diárias e histórico};
    \node[classe, below=of jparepo] (entity) {\textbf{Activity}\\distância, passos, velocidade,\\validade, risco e data};

    \draw[dep] (screen) -- (vm);
    \draw[dep] (vm) -- (repoi);
    \draw[impl] (repo) -- (repoi);
    \draw[dep] (vm) -- (sensor);
    \draw[dep] (sensor) -- (evaluator);
    \draw[dep] (repo) -- node[above, align=center]{REST/JSON} (controller);
    \draw[dep] (controller) -- (activityservice);
    \draw[dep] (activityservice) -- (fraud);
    \draw[dep] (activityservice) -- (game);
    \draw[dep] (activityservice) -- (jparepo);
    \draw[dep] (jparepo) -- (entity);
\end{tikzpicture}
}
    \caption{Diagrama de classes simplificado do registro de atividades}
    \label{fig:classes_fitworkup}
\end{figure}

No cliente, \texttt{WorkoutScreen} apresenta o estado mantido por \texttt{WorkoutViewModel}. O serviço de sensores coleta os dados e o avaliador local produz indicadores de risco. A implementação do repositório envia o registro para o controlador da API. No servidor, \texttt{ActivityService} coordena a validação, a persistência e a concessão de recompensas, delegando responsabilidades específicas aos serviços de detecção e gamificação.

\section{Modelagem do Banco de Dados}
\label{sec:fitworkup_banco}

O modelo remoto contém entidades para usuários, atividades, amizades, conquistas, sessões em grupo, itens da loja, inventário e efeitos temporários. A Figura~\ref{fig:der_fitworkup} apresenta os relacionamentos mais relevantes. Campos secundários foram omitidos para preservar a legibilidade.

\begin{figure}[H]
    \centering
    \resizebox{\textwidth}{!}{\begin{tikzpicture}[
    font=\scriptsize,
    node distance=0.65cm and 0.75cm,
    entidade/.style={draw, rounded corners, align=left, text width=3.15cm, inner sep=5pt, fill=blue!5},
    rel/.style={->, >=stealth, thick}
]
    \node[entidade] (user) {\textbf{users}\\\underline{id}, username, email, password, weight\_kg, xp, level, fitcoins, streak, avatar\_key, active};
    \node[entidade, right=of user] (activity) {\textbf{activities}\\\underline{id}, user\_id, type, distance\_km, duration\_seconds, steps, avg\_speed, is\_valid, risk\_score, timestamp};
    \node[entidade, below=of user] (friendship) {\textbf{tb\_friendships}\\\underline{id}, user\_id, friend\_id, status, created\_at};
    \node[entidade, below=of activity] (group) {\textbf{group\_sessions}\\\underline{id}, host\_user\_id, code, name, target\_distance\_km, status, created\_at};
    \node[entidade, right=of group] (participant) {\textbf{group\_participants}\\\underline{id}, group\_session\_id, user\_id, ready, joined\_at};

    \node[entidade, below=of friendship] (ua) {\textbf{tb\_user\_achievements}\\\underline{id}, user\_id, achievement\_id, unlocked\_at};
    \node[entidade, left=of ua] (achievement) {\textbf{tb\_achievements}\\\underline{id}, name, description, criteria\_type, target\_value, xp\_reward, fitcoins\_reward};
    \node[entidade, below=of group] (inventory) {\textbf{inventory\_items}\\\underline{id}, user\_id, store\_item\_id, quantity, is\_equipped};
    \node[entidade, right=of inventory] (store) {\textbf{store\_items}\\\underline{id}, name, price, category, effect\_type, multiplier, duration\_minutes, active};
    \node[entidade, below=of inventory] (boost) {\textbf{user\_boosts}\\\underline{id}, user\_id, store\_item\_id, effect\_type, multiplier, starts\_at, expires\_at};

    \draw[rel] (user) -- node[above]{1:N} (activity);
    \draw[rel] (user) -- node[left]{1:N} (friendship);
    \draw[rel] (user) -- node[above, sloped]{1:N} (group);
    \draw[rel] (group) -- node[above]{1:N} (participant);
    \draw[rel] (user) -- node[above, sloped]{1:N} (participant);
    \draw[rel] (user) -- node[above, sloped]{1:N} (ua);
    \draw[rel] (achievement) -- node[above]{1:N} (ua);
    \draw[rel] (user) -- node[above, sloped]{1:N} (inventory);
    \draw[rel] (store) -- node[above]{1:N} (inventory);
    \draw[rel] (user) -- node[left]{1:N} (boost);
    \draw[rel] (store) -- node[right]{1:N} (boost);
\end{tikzpicture}
}
    \caption{Diagrama entidade-relacionamento simplificado do FitWorkUp}
    \label{fig:der_fitworkup}
\end{figure}

A entidade \texttt{User} é central no modelo. Uma pessoa pode possuir várias atividades, amizades, conquistas, participações em grupos e entradas de inventário. \texttt{UserAchievement} resolve a associação entre usuários e conquistas; \texttt{InventoryItem} associa usuários e produtos; e \texttt{GroupParticipant} representa a participação em uma sessão coletiva. Os efeitos temporários adquiridos na loja são registrados separadamente em \texttt{UserBoost}.

\section{Interfaces e Fluxos do Aplicativo}
\label{sec:fitworkup_interfaces}

Esta seção deve utilizar capturas reais da versão avaliada do aplicativo. Para evitar divergência entre texto e implementação, recomenda-se capturar todas as telas no mesmo dispositivo, resolução e tema. Os blocos abaixo compilam mesmo antes da inclusão das imagens.

\providecommand{\telaFitWorkUp}[4]{%
\begin{figure}[H]
    \centering
    \IfFileExists{#1}{%
        \includegraphics[width=#2\textwidth,keepaspectratio]{#1}%
    }{%
        \fbox{\parbox[c][6cm][c]{0.72\textwidth}{\centering
        Inserir captura real em: \\ \texttt{#1}}}%
    }
    \caption{#3}
    \label{#4}
\end{figure}}

\subsection{Autenticação e Recuperação de Senha}

A tela de autenticação recebe as credenciais e direciona o usuário à área principal quando a API confirma a sessão. O fluxo de recuperação solicita e valida o código enviado ao endereço cadastrado antes de permitir a alteração da senha.

\telaFitWorkUp{Figuras/prints/login.png}{0.32}{Tela de autenticação}{fig:tela_login}

\subsection{Tela Inicial e Resumo de Progresso}

A tela inicial concentra o acesso aos módulos e apresenta dados resumidos do usuário. A partir dela, é possível iniciar uma atividade individual, criar ou entrar em uma sessão em grupo e consultar perfil, ranking e loja.

\telaFitWorkUp{Figuras/prints/inicio.png}{0.32}{Tela inicial do FitWorkUp}{fig:tela_inicio}

\subsection{Registro de Atividade}

Durante a atividade, a interface exibe as métricas registradas e o percurso. Os controles permitem iniciar, pausar, retomar e finalizar o treino. Ao concluir, os dados são encaminhados para avaliação e persistência.

\telaFitWorkUp{Figuras/prints/atividade.png}{0.32}{Acompanhamento de uma atividade}{fig:tela_atividade}

\subsection{Perfil, Ranking e Loja}

O perfil apresenta progressão e conquistas; o ranking organiza o desempenho semanal; e a loja permite consultar itens, saldo e inventário. Recomenda-se inserir uma captura de cada módulo ou uma composição com três imagens alinhadas.

\telaFitWorkUp{Figuras/prints/perfil.png}{0.32}{Perfil e progressão do usuário}{fig:tela_perfil}
\telaFitWorkUp{Figuras/prints/ranking.png}{0.32}{Ranking semanal}{fig:tela_ranking}
\telaFitWorkUp{Figuras/prints/loja.png}{0.32}{Loja virtual e itens disponíveis}{fig:tela_loja}

\section{Limitações da Implementação}
\label{sec:fitworkup_limitacoes}

Na versão examinada, o mecanismo de validação utiliza regras e limiares definidos no cliente e no servidor; seu desempenho precisa ser medido em cenários controlados antes de qualquer afirmação de eficácia. As métricas produzidas pelos sensores variam conforme dispositivo, posicionamento, qualidade do sinal e permissões do sistema. A estimativa de calorias não possui finalidade clínica.

Também foi identificada uma dependência residual do driver do SQL Server no arquivo de construção da API, embora o perfil de desenvolvimento esteja configurado para PostgreSQL. Essa dependência deve ser removida para evitar ambiguidade na documentação. Por fim, o serviço de mensagens de incentivo é baseado em regras; a integração com Google Gemini e \textit{Spring AI} não está implementada na versão analisada.

\section{Considerações Finais do Capítulo}
\label{sec:fitworkup_consideracoes}

O capítulo apresentou o FitWorkUp a partir de sua implementação, descrevendo funcionalidades, tecnologias, arquitetura, classes, banco de dados e interfaces. Essa documentação fornece a base para o capítulo de avaliação, no qual os requisitos e os componentes centrais devem ser confrontados com casos de teste e cenários controlados.
